\subsection{cycle}\label{cycle}
Das \verb|cycle| Modul bildet die zeitliche Grundstruktur des Slot\-basierten
Funkprotokolls ab. Ein vollständiger \emph{Cycle} wird in \verb|SLOT_CNT|
Slots unterteilt, und jeder Slot wiederum in \verb|SUB_SLOT_CNT| Sub\-Slots.
Mit den Standardwerten aus \verb|system_definitions.h| (Listing
\ref{cycle_defines}) ergibt das $16$ Slots zu je $8$ Sub\-Slots, also $128$
Sub\-Slots pro Cycle. Jedes Gerät besitzt einen eigenen, ungeraden Slot
($1, 3, 5, \dots, \mathtt{f}$) und darf nur in diesem senden; die übrigen
Slots dienen dem Empfang der anderen Teilnehmer.

\begin{listing}
\begin{minted}{C}
#define SUB_SLOT_POW2  3
#define SUB_SLOT_CNT   (1<<SUB_SLOT_POW2)   // 8 Sub-Slots pro Slot
#define SUB_SLOT_MASK  (SUB_SLOT_CNT-1)     // 0x07
#define SLOT_POW2      4
#define SLOT_CNT       (1<<SLOT_POW2)        // 16 Slots pro Cycle
#define SLOT_MASK      (SLOT_CNT-1)          // 0x0F
#define SLOT_SHIFT     (SUB_SLOT_POW2)       // 3
\end{minted}
\caption{Slot\-Konstanten aus \texttt{system\_definitions.h}}
\label{cycle_defines}
\end{listing}

\subsubsection*{Zustand}
Der gesamte Zustand wird in der Struktur \verb|cycle_t| (Listing
\ref{cycle_t}) gehalten. Das zentrale Feld ist der monoton hochzählende
Zähler \verb|subSlot| (Bereich $0 \dots 127$), in dem Slot\- und
Sub\-Slot\-Index \emph{bitweise gepackt} sind: die unteren
\verb|SUB_SLOT_POW2| Bits enthalten den Sub\-Slot, die darüber liegenden
\verb|SLOT_POW2| Bits den Slot. Die beiden Makros aus Listing
\ref{cycle_macros} extrahieren die beiden Anteile, ohne dass an anderer
Stelle gerechnet werden muss:
\begin{align*}
\mathtt{ACT\_SUB\_SLOT} &= \mathtt{subSlot}\ \&\ \mathtt{SUB\_SLOT\_MASK} \\
\mathtt{ACT\_SLOT}      &= (\mathtt{subSlot} \gg \mathtt{SLOT\_SHIFT})\ \&\ \mathtt{SLOT\_MASK}
\end{align*}
Die Felder \verb|actSlot| und \verb|sSlot| sind zwischengespeicherte Kopien
des aktuellen Slots bzw. Sub\-Slots, die bei jedem Fortschalten neu gesetzt
werden, damit sie ohne erneute Bitoperation zur Verfügung stehen.
\verb|cycle| zählt die vollständig durchlaufenen Cycles ($16$ Bit, läuft
bei $65535$ über) und \verb|init| markiert die erfolgte Initialisierung.

\begin{listing}
\begin{minted}{C}
typedef struct cycle_s {
    volatile int8_t subSlot; // aktueller SubSlot (actSlot+sSlot gepackt)
    int8_t          actSlot; // zwischengespeicherter Slot
    int8_t          sSlot;   // zwischengespeicherter Sub-Slot
    uint16_t        cycle;   // Anzahl durchlaufener Cycles
    bool            init;    // true nach cycle_init()
} cycle_t;
\end{minted}
\caption{Definition der \texttt{cycle\_t} Struktur}
\label{cycle_t}
\end{listing}

\begin{listing}
\begin{minted}{C}
#define ACT_SUB_SLOT(_cycle)  (((_cycle)->subSlot) & SUB_SLOT_MASK)
#define ACT_SLOT(_cycle)      (((_cycle)->subSlot >> SLOT_SHIFT) & SLOT_MASK)
\end{minted}
\caption{Makros zum Extrahieren von Slot und Sub\-Slot}
\label{cycle_macros}
\end{listing}

\subsubsection*{API}
\begin{description}
\item[\texttt{cycle\_init(cycle)}] Setzt \verb|init| und ruft anschliessend
  \verb|cycle_reset()| auf. Rückgabe: \verb|em_msg| (\verb|EM_OK|, bzw.\
  \verb|EM_ERR| bei einem Nullzeiger).
\item[\texttt{cycle\_reset(cycle)}] Setzt \verb|subSlot| und \verb|cycle|
  auf $0$ zurück. Erfordert eine zuvor initialisierte Struktur. Rückgabe:
  \verb|em_msg| (\verb|EM_ERR|, wenn nicht initialisiert).
\item[\texttt{cycle\_set\_slot(cycle, slot)}] Positioniert den Zähler auf den
  Beginn des angegebenen Slots ($\mathtt{subSlot} = \mathtt{slot} \cdot
  \mathtt{SUB\_SLOT\_CNT}$). Nur gültige Slots werden akzeptiert; der Slot wird
  zuvor mit \verb|cycle_check_slot()| geprüft. Rückgabe: \verb|em_msg|
  (\verb|EM_OK| bei gültigem Slot, \verb|EM_ERR| bei Nullzeiger oder
  ungültigem Slot).
\item[\texttt{cycle\_check\_slot(slot)}] Validiert einen Slot: zulässig sind
  ungerade Werte im Bereich $0 < \mathtt{slot} \le \mathtt{SLOT\_CNT}$.
  Rückgabe: \verb|int8_t| -- der Slot oder $-1$.
\item[\texttt{cycle\_check(cycle, rxSlot, ss)}] Prüft, ob ein vom Slot
  \verb|rxSlot| empfangenes Frame zur aktuellen Cycle\-Position passt. Das
  Ergebnis ist wahr, wenn der aktuelle Slot gleich \verb|rxSlot| ist, oder
  wenn \verb|rxSlot| der unmittelbar \emph{folgende} Slot ist und der aktuelle
  Sub\-Slot kleiner als die Toleranz \verb|ss| ist, oder wenn \verb|rxSlot|
  der unmittelbar \emph{vorangehende} Slot ist und weniger als \verb|ss|
  Sub\-Slots bis zum Slotende verbleiben (Slotindizes jeweils modulo
  \verb|SLOT_CNT|). Das Toleranzfenster \verb|ss| lässt damit auch Frames der
  direkten Nachbarslots zu. Rückgabe: \verb|bool| (\verb|false| bei Nullzeiger
  oder ungültigem \verb|rxSlot|).
\item[\texttt{cycle\_difference(cycle, rxSlot)}] Liefert den vorzeichen\-
  behafteten Abstand in Sub\-Slots von der aktuellen Position \verb|subSlot|
  zum Beginn des Slots \verb|rxSlot| ($\mathtt{rxSlot} \cdot
  \mathtt{SUB\_SLOT\_CNT}$), gemessen auf dem kürzesten Weg über den Ring aus
  $\mathtt{SUB\_SLOT\_CNT} \cdot \mathtt{SLOT\_CNT} = 128$ Sub\-Slots. Der
  Wert liegt im Bereich $[-64, 63]$; negativ bedeutet, dass \verb|rxSlot|
  voraus liegt, positiv, dass er bereits zurückliegt:
  \[
    \mathtt{diff} = \big((\mathtt{subSlot} - \mathtt{rxSlot} \cdot
    \mathtt{SUB\_SLOT\_CNT} + 64) \bmod 128\big) - 64 .
  \]
  Rückgabe: \verb|int8_t| ($0$ bei Nullzeiger oder nicht initialisierter
  Struktur).
\item[\texttt{cycle\_string(cycle)}] Liefert eine kompakte Darstellung im
  Format \verb|(c:<cycle>, <slot>, <subSlot>)| in einem statischen Puffer.
  Rückgabe: \verb|char *| (\verb|NULL| bei Nullzeiger).
\item[\texttt{cycle\_print(cycle, title)}] Gibt alle Felder zeilenweise aus
  (Debug\-Hilfe, vgl. \verb|xx_print| Konvention in der Einführung). Rückgabe:
  \verb|em_msg| (\verb|EM_ERR| bei Nullzeiger oder nicht initialisierter
  Struktur, sonst \verb|EM_OK|).
\end{description}

\subsubsection*{Fortschalten des Cycles}
{\sloppy\emergencystretch=3em
Der Zähler wird ausschliesslich über \verb|cycle_increment()| weitergeschaltet,
typischerweise einmal pro Sub\-Slot\-Tick. Die Funktion ist an die
Synchronisations\-Zustandsmaschine (\verb|system_state_e|) gekoppelt:

\begin{itemize}
\item Im Zustand \verb|SYNCHRONIZE| wird der Übergang nach
  \verb|SYNCHRONIZE_READY| ausgelöst und das interne Flag \verb|is_set|
  gesetzt; ab diesem Zeitpunkt läuft der Zähler.
\item Solange \verb|is_set| gilt, wird \verb|subSlot| inkrementiert und modulo
  $\mathtt{SUB\_SLOT\_CNT} \cdot \mathtt{SLOT\_CNT}$ ($=128$) zurückgefaltet;
  \verb|actSlot| und \verb|sSlot| werden aus dem neuen Wert nachgeführt.
\item In den Zuständen \verb|SYNCHRONIZE_DOING|, \verb|SYNCHRONIZE_READY| und
  \verb|SYNCHRONIZE_ERROR| wird ein Slotwechsel erkannt. Beim Übergang auf
  Slot $0$ wird der Cycle\-Zähler \verb|cycle| genau einmal pro Umlauf erhöht
  (über das Flag \verb|cycle_once| entprellt).
\end{itemize}
\par}

\noindent Da \verb|cycle_increment()| zuerst fortschaltet und erst danach die
Felder meldet, liefert der erste Aufruf nach einem Reset bereits den nächsten
Sub\-Slot. Ist \verb|OPTION_SHOW_TIMING| aktiv, werden zusätzlich Status\-LEDs
zur Laufzeitmessung getoggelt.
